\section{The full-turn cap problem}\label{sec:full-turn}

We first determine the geometry of a maximizing cap. The support
functional in the next section will then identify its value and its
shape. Throughout this section $L=\pi/2$ and $P=P_L$.

\subsection{The two charged wings}

Let $U$ be a prescribed canonical full maximum. Translate its
projection to $[-2C,2C]$, so $J_U=[-C,C]$. If the affine middle
roof has positive slope, put
\[
 f_U(-C)=1-h,\qquad f_U(C)=1,\qquad
 \operatorname{top}(U)=[C,C+T]\times\{1\}.
\]
The higher endpoint is one by \cref{lem:selection-top}, and the
middle facet ends at $C$ by \cref{lem:variation-middle-junctions}.
The lower endpoint is greater than $1/2$ by
\cref{lem:variation-endpoint-pressure}. Thus $0<h<1/2$.
The width and height estimates of \cref{lem:preliminary-cuts},
followed by \cref{cor:full-three-tent-cut}, sharpen this to
\begin{equation}\label{eq:tilted-parameter-range}
 1/2<C<37/50,\qquad 0<h<17/50.
\end{equation}
The strict width cut applies with $a=2C$: a cap with
$C\ge37/50$ would have $P<41/50$, whereas a maximum has
$P\ge\Mc/2>41/50$.

We use the actual support of the high wing and the support of the
low wing alone:
\begin{equation}\label{eq:tilted-wing-supports}
 F(t)=h_U(t),\qquad
 G(t)=\max_{-2C\le x\le-C}
                 \{-x\sin t+f_U(x)\cos t\},\quad 0\le t\le L.
\end{equation}
In this section $G$ omits the middle facet. Its endpoint values
and derivatives are
\begin{equation}\label{eq:tilted-wing-traces}
 \begin{gathered}
 F(0)=G(L)=2C,\quad F(L)=1,\quad G(0)=1-h,\\
 F'(0)=e_R,\quad F'(L)=-C-T,\quad
 G'(0)=C,\quad G'(L)=-e_L,
 \end{gathered}
\end{equation}
where $e_R=f_U(2C)$ and $e_L=f_U(-2C)$.
Let
\begin{equation}\label{eq:tilted-contact-data}
 \begin{gathered}
 u=F''+F,\quad v=G''+G,\qquad p=F'-G+1,\quad q=G'+F-1,\\
 z=(F-1)u(t)+(G-1)u^\perp(t),\\
 B=z+p u^\perp(t),\qquad D=z-q u(t).
 \end{gathered}
\end{equation}
Here $u$ and $v$ on the first line are scalar curvature densities;
$u(t)$ and $u^\perp(t)$ on the second line are the normal vectors
defined earlier. To distinguish the niche endpoint value from the
corner vector, below we write $z_+=N_L(U)(C)$.
Differentiation gives, almost everywhere,
\begin{equation}\label{eq:tilted-contact-kinematics}
 p'=u-1-q,\quad q'=v-1+p,\qquad
 B'=(u-1)u^\perp(t),\quad D'=(1-v)u(t).
\end{equation}
The tangent endpoints are
\begin{equation}\label{eq:tilted-tangent-endpoints}
 \begin{gathered}
 B(0)=(2C-1,e_R),\qquad B(L)=(C+T,0),\\
 D(0)=(-C,-h),\qquad D(L)=(1-2C,e_L).
 \end{gathered}
\end{equation}

\begin{lemma}[Removing the middle facet]\label{lem:tilted-surrogate}
The supports in \cref{eq:tilted-wing-supports} belong to
$W^{2,\infty}(0,L)\cap C^1[0,L]$. Their curvature densities are
the charged wing densities of \cref{thm:variation-curvature}.
The positive niche on $J_U$ is unchanged if the actual companion
support $h_U(t+L)$ is replaced by $G(t)$.
For $0<t<\delta:=\arctan(h/(2C))$ one has $u=v=0$.
\end{lemma}
\begin{proof}
Concavity of the roof and the pinned middle facet give
\[
 h_U(t+L)=\max\{G(t),-C\sin t+\cos t\}.
\]
The added support term has inner-wall height
$1-\sec t+(x-C)\tan t\le0$ for $x\in J_U$.
Distribute the two-wall minimum over this maximum and take the
positive part; this proves the niche assertion angle by angle.
The support $G$ starts with the point $(-C,1-h)$ and contains no
middle-facet atom. For $t\le\delta$ that point dominates every
point of the low wing, by the supporting middle line. The actual
companion is instead attained at $(C,1)$, whose wall is
nonpositive on $J_U$. Thus no positive inner graph can carry
either source at these angles. The exact source identity of
\cref{lem:variation-sources} gives $u=v=0$ there. For
$t>\delta$, the low wing gives the actual companion support.
The regularity and density bounds of
\cref{thm:variation-curvature} therefore apply to $G$; matching
the zero-density initial interval introduces no atom, because
the middle-junction inequalities are strict. The traces are the
coordinates of the extremal support points, giving
\cref{eq:tilted-wing-traces}.
\end{proof}

\subsection{Excluding an inclined middle}

The essential estimate is that a first wing with curvature at
most one cannot support an inclined maximum. The companion wing
may have curvature greater than one and may have tangencies
below the floor. The proof uses one joint source measure to keep
track of the fraction of a corner assigned to each wall.

\begin{proposition}[First-wing exclusion]
\label{lem:tilted-first-unit-exclusion}
A canonical full maximum with positive middle slope and
$1/2<C<4/5$, $0<h<1/2$ cannot have $u\le1$ almost everywhere.
\end{proposition}
\begin{proof}
The contact geometry and the joint source decomposition are proved
in \cref{lem:tilted-first-unit-geometry,lem:variation-joint-occupation}.
If the companion density is also at most one, the two clipped
strip moments give the contradiction
$6CT\le4T/3<3T$ in \cref{lem:tilted-two-unit}.
Otherwise \cref{lem:tilted-fold-energy} proves
\[
 T>0,\qquad 6CT\le\frac{97}{48}T<3T,
\]
contradicting $C>1/2$. The fractional source equations used there
allow measurable ties, folds, and later floor crossings.
\end{proof}

\begin{lemma}[Horizontal middle rigidity]\label{lem:full-affine-rigidity}
Every canonical full maximum has middle slope zero.
\end{lemma}
\begin{proof}
Suppose first that its middle slope is positive. The parameter
range is \cref{eq:tilted-parameter-range}. For $C\le2/3$, or
for $h\ge21/100$, \cref{lem:tilted-unit-regions} gives $u\le1$,
contradicting \cref{lem:tilted-first-unit-exclusion}.

It remains that $2/3<C<37/50$ and $h<21/100$.
The endpoint geometry in \cref{lem:tilted-dry-endpoint} gives
$N_L(U)(-C)=0$. If $h\ge1/20$ and $C\le73/100$, the last
case of \cref{lem:tilted-unit-regions} again gives $u\le1$.
If $h\ge1/20$ and $C\ge73/100$, the improved three-angle
estimate \cref{lem:tilted-medium-width} gives $P<41/50$,
contradicting maximality.

The remaining interval is $0<h<1/20<509/10000$.
By \cref{lem:tilted-small-height}, failure of either initial-floor
energy test contradicts the ordinary one-turn area bound.
When both tests hold, \cref{lem:tilted-initial-energy} gives
$u\le1$ and the first-wing exclusion applies. These alternatives
exhaust all positive slopes.
Horizontal reflection preserves the full-angle cap score and
changes the sign of the middle slope. Apply the positive-slope
conclusion to the reflection to exclude negative slopes.
\end{proof}

The horizontal estimates of \cref{lem:full-horizontal-unit} now
place every canonical full maximum in the class with both wing
curvatures at most one, top face exactly $J_U$, and ordered
contact switches. We next evaluate the resulting support functional
and recover every original equality cap.
